\documentclass[11pt]{amsart}
\usepackage[T1]{fontenc}
\usepackage{lmodern}
\usepackage[margin=1in]{geometry}
\usepackage{amsmath,amssymb,amsthm}
\usepackage{microtype}
\usepackage[hidelinks]{hyperref}\hypersetup{pdftitle={Schur (2,2) reconstruction for explicit length classes},pdfauthor={Xiaojian Zeng},pdfsubject={Partial reconstruction theorem for the Schur (2,2) monomial-value multiset},pdfkeywords={Schur polynomial, multiset reconstruction, integer partition, logarithmic moment}}
\newtheorem{theorem}{Theorem}
\newtheorem{corollary}[theorem]{Corollary}
\newtheorem{lemma}[theorem]{Lemma}
\theoremstyle{remark}
\newtheorem{remark}[theorem]{Remark}
\newcommand{\Sxx}{\mathcal{S}_{22}}
\newcommand{\RR}{\mathbb{R}}
% Exceptional residues are written explicitly for citation-audit parsing.
% Multiset braces are written explicitly for citation-audit parsing.
\DeclareMathOperator{\SSYT}{SSYT}
\title[Schur $(2,2)$ reconstruction by length]{Schur $(2,2)$ reconstruction for explicit length classes}
\author{Xiaojian Zeng}
\address{PBC School of Finance, Tsinghua University, 100190, China}
\email{xiaojianzeng@foxmail.com}
\date{September 12, 2026}
\thanks{2020 Mathematics Subject Classification. Primary 05E05; Secondary 05A17.}
\keywords{Schur polynomial, multiset reconstruction, integer partition,
logarithmic moment, Chern shadow}

\begin{document}
\begin{abstract}
For a positive multiset, retain the value of every monomial of
$s_{(2,2)}$, including its coefficient as a multiplicity. We prove
reconstruction of every $n$-element positive real multiset whenever
$c_j(n)=n(n^2-1+3\cdot2^{j-1}-3^j)/3$ is nonzero for $1\leq j\leq n$.
The output cardinality and first $n$ logarithmic power sums suffice.
An elementary congruence argument verifies this condition for $n\geq3$,
$n\ne4$, outside the residue classes $14,26,34,46$ modulo $60$.
Separate arguments at lengths $2$ and $4$ give fixed-sum reconstruction.
Consequently the same-sum integer-partition question for shape $(2,2)$
holds at every length outside those four residue classes. The all-length
question, including length $14$, remains unresolved here. The contribution
is an application and arithmetic refinement of established moment and
Chern-shadow methods.
\end{abstract}
\maketitle

\section{The reconstruction problem and the result}

Let $a=\{\!\{a_1,\ldots,a_n\}\!\}$ be a multiset of positive real numbers,
where $n\geq2$. Indices distinguish occurrences, even when values coincide.
For each semistandard Young tableau $T$ of shape $(2,2)$ with entries
in $\{1,\ldots,n\}$, let $\alpha_i(T)$ be the number of occurrences of $i$.
Rows are weakly increasing and columns are strictly increasing. Define
\begin{equation}\label{eq:definition}
 \Sxx(a)=\{\!\{\prod_{i=1}^n a_i^{\alpha_i(T)}:
                    T\in\SSYT((2,2),n)\}\!\}.
\end{equation}
This is the multiset of monomial values in $s_{(2,2)}(a)$, counted with
the coefficients of the monomials. Equal numerical values contribute
additively to their multiplicity. Equality of these multisets is stronger
than equality of their supports or their sums.

An integer partition is identified with its multiset of positive integer
parts, written in nonincreasing order. The Schur partition maps and the
broader injectivity question appear in
Hadelyn--Niergarth--Li--Li \cite[Section~5, Question~5.3]{HadelynEtAl2026}.
Thomas--Tung \cite[Conjecture~6.3]{ThomasTung2026} ask whether every
two-row shape gives an injective map on the partitions of each prescribed
integer. For shape $(2,2)$ this asks, for every $N\geq2$, whether
\begin{equation}\label{eq:original}
 \lambda,\mu\text{ partitions of }N,\qquad \ell(\lambda),\ell(\mu)\geq2,
 \qquad \Sxx(\lambda)=\Sxx(\mu)
 \quad\Longrightarrow\quad \lambda=\mu.
\end{equation}
We prove this implication for explicit length classes. Its unrestricted
form remains unresolved by the results below.

For $j\geq1$ put
\begin{equation}\label{eq:coefficient}
 c_j(n)=\frac n3\bigl(n^2-1+3\cdot2^{j-1}-3^j\bigr),
 \qquad
 Q_j(a)=\sum_{v\in\Sxx(a)}(\log v)^j,
\end{equation}
with every sum over a multiset taken with multiplicities.

\begin{theorem}[Moment criterion]\label{thm:criterion}
Let $n\geq2$. If $c_j(n)\ne0$ for every $1\leq j\leq n$, then every
$n$-element positive real multiset $a$ is uniquely determined by
$|\Sxx(a)|$ and $Q_1(a),\ldots,Q_n(a)$. In particular $\Sxx$ is
injective on these multisets, allowing arbitrary repeated values and
without prescribing their sum.
\end{theorem}

\begin{theorem}[Sufficient length classes]\label{thm:classes}
If $n\geq3$, $n\ne4$, and $n\bmod60\notin\{14,26,34,46\}$, then $c_j(n)\ne0$
for every $j\geq1$. In particular Theorem~\ref{thm:criterion} applies
at every odd length $n\geq3$ and every length divisible by four with
$n\geq8$.
\end{theorem}

\begin{corollary}[Fixed-sum partitions]\label{cor:partitions}
Let $N\geq2$ be an integer, and let $\lambda,\mu$ be partitions of $N$ with at
least two parts. If $\Sxx(\lambda)=\Sxx(\mu)$, their lengths are equal.
If this common length is outside $\{14,26,34,46\}$ modulo $60$, then
$\lambda=\mu$. Lengths $2$ and $4$ are included.
\end{corollary}

These are sufficient conditions. Membership in an excluded residue class
does not imply failure of the coefficient criterion, and a zero coefficient
does not imply a collision of output multisets.

Thomas--Tung prove the $(2,1)$ case and show, for a shape with $k$ rows,
that prescribing the first $k-1$ parts suffices for recovery
\cite[Theorems~1.7 and~5.7]{ThomasTung2026}. For $(2,2)$, our length
criterion removes this need to know the largest part.
The underlying triangular power-sum argument is familiar from subset-sum
reconstruction; see Fomin \cite[equation~(1), Theorem~2.1]{Fomin2019}.
The coefficient in \eqref{eq:coefficient} is a specialization of the
Chern shadows of D'Andrea--Fatighenti--Onorati
\cite[Example~4.7, Appendix~A]{DAndreaEtAl2025}.
Our contribution is their application to the full $(2,2)$ tableau multiset
and the explicit arithmetic length criterion. The comparisons made here
identify the scope of this application; they do not assert comprehensive
priority.

Repeated values are a substantive part of a symmetric reconstruction
problem. In the related setting of trees, Zeng
\cite[Theorem~1.1, Lemmas~2.2 and~3.1]{Zeng2026MultiplicityArxiv}
recovers a proper tree from its chromatic symmetric function when every
nonleaf vertex of its weighted core has a weight occurring only once;
weights on core leaves may repeat. Here all input multiplicities are
allowed: they are recovered as root multiplicities by Newton identities.

\section{Tableau contents and logarithmic moments}

Write $m_\nu$ for the monomial symmetric polynomial of exponent partition
$\nu$. A label occurs at most twice in a tableau of shape $(2,2)$.
There is one tableau for each content of type $(2,2)$ or $(2,1,1)$,
and two for each content of type $(1,1,1,1)$. Indeed, for three ordered
labels $i<j<k$, the tableaux, displayed by rows, are
$(i,i;j,k)$, $(i,j;j,k)$, and $(i,j;k,k)$. For four ordered labels
$i<j<k<l$, they are $(i,j;k,l)$ and $(i,k;j,l)$. Thus
\begin{equation}\label{eq:contents}
 s_{(2,2)}=m_{22}+m_{211}+2m_{1111}.
\end{equation}
The output cardinality is therefore
\begin{equation}\label{eq:cardinality}
 D(n)=\binom n2+3\binom n3+2\binom n4
     =\frac{n^2(n^2-1)}{12},
\end{equation}
where a binomial coefficient is zero if its lower index exceeds $n$.
The first summand is strictly increasing for integers $n\geq2$ and the
others are nondecreasing, so $D(n)$ determines $n$.

Let $u_r=\sum_i x_i^r$. Expanding by exponent types gives
\begin{equation}\label{eq:character}
 s_{(2,2)}(x)=\frac{u_1^4+3u_2^2-4u_1u_3}{12}.
\end{equation}
For a direct check, the coefficients of
$(m_4,m_{31},m_{22},m_{211},m_{1111})$ in $u_1^4$ are
$(1,4,6,12,24)$, in $u_2^2$ they are $(1,0,2,0,0)$, and in
$u_1u_3$ they are $(1,1,0,0,0)$. The combination in
\eqref{eq:character} is exactly \eqref{eq:contents}.

\begin{proof}[Proof of Theorem~\ref{thm:criterion}]
Set $y_i=\log a_i$, $p_0=n$, and $p_j=\sum_i y_i^j$ for $j\geq1$.
Introduce the formal series
\[
 P(t)=\sum_i e^{t y_i}=\sum_{j\geq0}p_j\frac{t^j}{j!}.
\]
Substituting $x_i=e^{t y_i}$ in \eqref{eq:character} yields
\begin{equation}\label{eq:generating}
 \sum_{v\in\Sxx(a)}e^{t\log v}
   =\frac{P(t)^4+3P(2t)^2-4P(t)P(3t)}{12}.
\end{equation}
At total degree $j$, a term involving $p_j$ has all other factors of
degree zero. Its coefficient in the $j$th derivative at zero is
\[
 \frac{4n^3+6n\,2^j-4n(1+3^j)}{12}=c_j(n).
\]
Every other term uses only $p_1,\ldots,p_{j-1}$. Consequently
\begin{equation}\label{eq:triangular}
 Q_j(a)=c_j(n)p_j+R_{j,n}(p_1,\ldots,p_{j-1}),
\end{equation}
where $R_{j,n}$ is explicit: set $p_j=0$ in the degree-$j$ part of
the right-hand side of \eqref{eq:generating} and multiply by $j!$.
There is no analytic approximation in this identity.

By \eqref{eq:cardinality}, the data determine $n$. Since the first $n$
coefficients $c_j(n)$ are nonzero, \eqref{eq:triangular} successively
determines $p_1,\ldots,p_n$. Put $E_0=1$ and use Newton's recurrence
\begin{equation}\label{eq:newton}
 kE_k=\sum_{j=1}^k(-1)^{j-1}E_{k-j}p_j,
 \qquad 1\leq k\leq n.
\end{equation}
Then $E_k$ is the $k$th elementary symmetric function of the $y_i$, and
\[
 z^n-E_1z^{n-1}+\cdots+(-1)^nE_n=\prod_{i=1}^n(z-y_i).
\]
The polynomial determines its root multiset, including repeated roots.
Exponentiation recovers $a$. The data are assumed to arise from a
positive multiset, so no separate existence assertion for arbitrary
moment sequences is needed.
\end{proof}

For the attribution of the coefficient, write the derivation of
\cite[Definition~3.2]{DAndreaEtAl2025} as
$\mathsf D_s=\sum_{r\geq1}r^{s+1}\partial/\partial u_r$.
Equation~\eqref{eq:character} gives
\[
 \mathsf D_s(s_{(2,2)})
 =\frac{u_1^3-u_3}{3}+u_2\,2^s-u_1\,3^s.
\]
Rank evaluation $u_r=n$ and $s=j-1$ gives $c_j(n)$. This is the
Chern-shadow specialization described in that source. Our proof keeps
the lower-moment polynomial $R_{j,n}$ at every order; it does not extend
a low-degree scalar identity by assumption.

\section{The arithmetic condition}

\begin{proof}[Proof of Theorem~\ref{thm:classes}]
For $n>0$, the equation $c_j(n)=0$ is equivalent to
\begin{equation}\label{eq:square}
 n^2=A_j:=1+3^j-3\cdot2^{j-1}.
\end{equation}
For $j=1,2,3$, the values are $1,4,16$, giving $n=1,2,4$.
If $j\geq4$ is even, then $A_j\equiv2\pmod8$, which is not a square.
If $j\geq5$ is odd, then $A_j\equiv4\pmod8$, so a square root
must satisfy $n\equiv2\pmod4$. Also $A_j\equiv1\pmod3$, hence
$3\nmid n$.

For odd $j$, one has $A_j\equiv1\pmod5$: when $j\equiv1\pmod4$,
the pair $(3^j,2^{j-1})$ is $(3,1)$ modulo $5$; when
$j\equiv3\pmod4$, it is $(2,4)$. Thus $n\equiv\pm1\pmod5$.
The simultaneous necessary congruences are
\[
 n\equiv2\pmod4,\qquad n\equiv\pm1\pmod3,
 \qquad n\equiv\pm1\pmod5.
\]
Their four solutions modulo $60$ are $14,26,34,46$. Outside these
classes, with $n\geq3$ and $n\ne4$, \eqref{eq:square} has no solution
at any order $j\geq1$.
\end{proof}

\section{The fixed-sum boundary lengths}

For a positive multiset $a$, let $\mathcal E_k(a)$ consist of the products
over $k$-element subsets of its index set. For a positive multiset
$B=\{\!\{b_1,\ldots,b_m\}\!\}$, put
\[
 \mathcal H_2(B)=\{\!\{b_i b_j:1\leq i\leq j\leq m\}\!\}.
\]
The next elementary recovery argument is the complete homogeneous
case of the extremal reconstruction in
\cite[Theorem~5.1]{HadelynEtAl2026}; we give it for positive reals.

\begin{lemma}\label{lem:pair}
The map $\mathcal H_2$ is injective on nonempty positive multisets, and
\begin{equation}\label{eq:pairs}
 \mathcal H_2(\mathcal E_2(a))=\Sxx(a)\mathbin{\dot{\cup}}\mathcal E_4(a),
\end{equation}
where $\mathbin{\dot{\cup}}$ adds multiplicities.
\end{lemma}

\begin{proof}
Order $b_1\geq\cdots\geq b_m>0$. The cardinality $m(m+1)/2$
determines $m$, and the maximum of $\mathcal H_2(B)$ is $b_1^2$.
After recovering $b_1,\ldots,b_r$, subtract all products from this
prefix, with their multiplicities. If $r<m$, the maximum residual is
$b_1b_{r+1}$: every surviving pair contains an index outside the prefix,
and its other factor is at most $b_1$. This determines the next value.
The argument permits repeated maxima.

For \eqref{eq:pairs}, classify two two-element index subsets. Coincident
pairs yield exponent type $(2,2)$ once; pairs meeting in one index yield
each type $(2,1,1)$ once; disjoint pairs yield three copies for each
four-element subset. Formula~\eqref{eq:contents} gives only two copies
in the last case, leaving exactly $\mathcal E_4(a)$.
\end{proof}

\begin{lemma}\label{lem:boundary}
At lengths $2$ and $4$, $\Sxx(a)$ and the prescribed sum
$\sum_i a_i=N>0$ determine every positive real multiset $a$.
\end{lemma}

\begin{proof}
At length two the only output value is $(a_1a_2)^2$. Positivity gives
the product, and the sum gives the quadratic with roots $a_1,a_2$.

At length four, put $A=a_1a_2a_3a_4$. There are $D(4)=20$ output
entries, each of degree four. Symmetry makes their total exponent the
same in every input variable, so
\[
 \prod_{v\in\Sxx(a)}v=A^{20}.
\]
The output therefore determines the positive number $A$. Since
$\mathcal E_4(a)=\{\!\{A\}\!\}$, adding one copy of $A$ in
\eqref{eq:pairs} and applying Lemma~\ref{lem:pair} recovers
$\mathcal E_2(a)$. Its sum is the elementary symmetric value $e_2(a)$.
The identity
\begin{equation}\label{eq:elementary}
 s_{(2,2)}(a)=e_2(a)^2-e_1(a)e_3(a)
\end{equation}
follows by exponent types: $e_2^2$ has coefficients $1,2,6$ at
$m_{22},m_{211},m_{1111}$, while $e_1e_3$ has $0,1,4$.
The left side is the sum of the observed entries. Since $e_1(a)=N>0$,
this determines $e_3(a)$. Together with $e_4(a)=A$, these data determine
\[
 z^4-Nz^3+e_2(a)z^2-e_3(a)z+A=\prod_{i=1}^4(z-a_i),
\]
and hence its positive root multiset.
\end{proof}

Li's fixed-sum pair-product theorem
\cite[Theorem~1 and the following positive-real extension]{Li2026}
also gives the final recovery once $\mathcal E_2(a)$ is known. The
boundary reductions are included to make the partition consequence
self-contained; they are not separate novelty claims.

\begin{proof}[Proof of Corollary~\ref{cor:partitions}]
Equation~\eqref{eq:cardinality} forces equal input lengths. At lengths
$2$ and $4$, apply Lemma~\ref{lem:boundary}. At every other length in
the stated classes, apply Theorems~\ref{thm:criterion} and~\ref{thm:classes}.
\end{proof}

The fixed-sum condition cannot simply be dropped from the boundary
statement. The inputs $(4,1)$ and $(2,2)$ both give $\{\!\{16\}\!\}$, but
have sums $5$ and $4$. At length four, $(8,2,2,2)$ and $(4,4,4,1)$
both give six copies of $16$, eight of $64$, and six of $256$, but
have sums $14$ and $13$. These examples do not contradict
\eqref{eq:original}.

\section{Exact checks and remaining scope}

The proofs above are symbolic. Supplementary integer and rational
calculations check the tableau multiplicities, the series identity and
Newton recovery without approximating logarithms. Taking $a_i=2^{y_i}$
allows rational $y_i$ to serve as exact base-two logarithmic coordinates;
the change of logarithm base rescales each moment by a known power.
One implementation checked all $1281$ nondecreasing lists over
$\{-2,-1,0,1,2\}$ of lengths $2$ through $8$, and $60$ specified
structured lists at lengths up to $61$. Of the $1341$ moment recoveries,
$1252$ complete and $89$ stop at the predicted coefficient zero.
Separate fixed-sum recovery succeeds for all $85$ small boundary inputs.
A second implementation, using a different log alphabet, checked $1281$
lists and $85$ rational boundary inputs. These are finite implementation
checks, not an exhaustive collision search over positive real inputs.

The accompanying coefficient certificate also gives a finite observation:
for $2\leq n\leq10^6$, the only zero-coefficient pairs at any order are
$(n,j)=(2,2),(4,3),(14,5)$. Indeed, exact integer square-root checks
of $A_j$ for $1\leq j\leq26$ give squares only at $j=1,2,3,5$,
with roots $1,2,4,14$. The stopping argument is
\[
 A_{j+1}-A_j=2\cdot3^j-3\cdot2^{j-1}>0\quad(j\geq1),
 \qquad A_{26}=2541765165034>10^{12}.
\]
This is a bounded classification and gives no classification at arbitrary
lengths. A short reproducer is included in the accompanying source package.

Length $14$ has $c_5(14)=0$, so the particular successive division in
\eqref{eq:triangular} does not recover $p_5$ at that step. Higher moments
or the prescribed input sum might supply additional information. No
noninjectivity statement at length $14$, converse to the criterion,
numerical stability bound, or all-length resolution is claimed.

\subsection*{Computational and AI assistance}
OpenAI Codex assisted with the moment calculation, proof development,
literature comparison, Python verification programs and manuscript
preparation. A separate model context reconstructed the argument and
wrote a second exact-arithmetic implementation. These checks are not
proof-assistant certification or journal peer review. The mathematical
claims rest on the proofs displayed above. The accompanying source
package documents the executed checks and their finite ranges.

\bibliographystyle{amsplain}
\bibliography{references}
\end{document}
